\section{Zadanie 5}
\autor{Kacper Orszulak}
Z kwadratu o boku 1 losujemy jednostajnie punkt \(P\). Punkt \(P\) wraz z bokami kwadratu tworzy 4 trójkąty. Wyznacz oczekiwane pole największego trójkąta.

\subsection*{Rozwiązanie}
Niech \(P = (X, Y)\), gdzie \(X\) i \(Y\) są niezależnymi zmiennymi losowymi z rozkładem jednostajnym na przedziale \([0, 1]\).
\begin{align*}
	\ev{A_{max}} & = \ev{max(A_1, A_2, A_3, A_4)}                                                \\
	             & = \ev{max\left(\frac{X}{2}, \frac{1-X}{2}, \frac{Y}{2}, \frac{1-Y}{2}\right)} \\
	             & = \frac{1}{2} \ev{max\left(X, 1-X, Y, 1-Y\right)}                             \\
	             & = \frac{1}{2} \ev{\underbrace{max(max(X, 1-X), max(Y, 1-Y))}_Z}               \\
	             & = \frac{1}{2} \int_{-\infty}^{\infty} t \cdot f_Z(t)dt                        \\
	             & = \frac{1}{2} \int_{\frac{1}{2}}^{1} t \cdot f_Z(t)dt
	\stepcounter{equation}\tag{\theequation}\label{myeq1}
\end{align*}
\begin{align*}
	F_Z(z) & = \mathcal{P}(Z \le z)                                                                                                                             \\
	       & = \mathcal{P}(\underset{\nwarrow \text{ Zmienne niezależne } \nearrow}{\underbrace{max(X, 1-X)} \le z \; \land \; \underbrace{max(Y, 1-Y)}} \le z) \\
	       & = \mathcal{P}(max(X, 1-X) \le z) \cdot \mathcal{P}(max(Y, 1-Y) \le z)                                                                              \\
	       & = F_{max(X, 1-X)}(z) \cdot F_{max(Y, 1-Y)}(z)
	\stepcounter{equation}\tag{\theequation}\label{myeq2}
\end{align*}
\begin{align*}
	F_{max(X, 1-X)}(z)                    & = F_{\frac{1}{2} + \left| X - \frac{1}{2} \right|}(z) = F_{\left| X - \frac{1}{2} \right|}\left(z - \frac{1}{2}\right) \\
	F_{\left| X - \frac{1}{2} \right|}(z) & = \mathcal{P}\left(\left|X-\frac{1}{2}\right| \le z\right)                                                             \\
	                                      & = \mathcal{P}\left(\frac{1}{2} - z \le X \le \frac{1}{2} + z\right)                                                    \\
	                                      & = \int_{\frac{1}{2} - z}^{\frac{1}{2} + z}f_X(t)dt = 2z
\end{align*}
Oczywiście zmienne \(X\) oraz \(Y\) mają ten sam rozkład, więc podstawiamy wynik pod ~\eqref{myeq2}.
\begin{align*}
	F_Z(z) & = \left(2z-1\right)^2                       \\
	f_Z(z) & = F_Z'(z) = 2(2z-1)(2z-1)' = 4(2z-1) = 8z-4
\end{align*}
Podstawiamy teraz pod ~\eqref{myeq1}.
\begin{align*}
	\ev{A_{max}} = \frac{1}{2} \int_{\frac{1}{2}}^{1} t \cdot f_Z(t)dt = \frac{1}{2}  \int_{\frac{1}{2}}^1 t \cdot (8t-4)dt = \frac{1}{2} \int_{\frac{1}{2}}^1 (8t^2-4t) dt = \frac{5}{12}
\end{align*}

\subsection*{Odpowiedź}
\[
	\boxed{\frac{5}{12}}
\]
